\conclusion

В ходе выполнения организационно-правового анализа были рассмотрены
основные нормативные правовые акты и нормативные документы, значимые для
разработки методики верификации соответствия TLA+ спецификации и исходного
кода. Проведённый анализ показал, что рассматриваемая методика относится к
сфере информационных технологий, обработки и защиты информации, а также
связана с современными подходами к безопасной разработке программного
обеспечения \cite{constitution_rf,fz149_info,doctrine_ib_rf,gost_56939_2024,fstek_240_2023}.

Конституция Российской Федерации формирует базовые правовые принципы,
которые должны учитываться при разработке и применении методики:
законность обращения с информацией, соблюдение прав граждан на
неприкосновенность частной жизни и допустимость ограничений доступа к
отдельным категориям сведений только на основании закона
\cite{constitution_rf}. Федеральный закон №149-ФЗ, в свою очередь,
определяет общую правовую рамку для работы с информацией, применения
информационных технологий и защиты информации, что непосредственно
соотносится с процессами трассировки, хранения, агрегации и анализа данных
исполнения распределённой системы \cite{fz149_info}.

Доктрина информационной безопасности Российской Федерации подчёркивает
значимость устойчивого функционирования информационной инфраструктуры,
развития безопасных информационных технологий и повышения надёжности
программных средств \cite{doctrine_ib_rf}. ГОСТ Р 56939--2024 и приказ
ФСТЭК России №240 дополняют данную правовую основу с точки зрения
безопасной разработки программного обеспечения, закрепляя требования к
организации соответствующих процессов и подтверждая значимость методов,
направленных на своевременное выявление недостатков и повышение качества
программных средств \cite{gost_56939_2024,fstek_240_2023}.

Таким образом, предлагаемая методика верификации соответствия TLA+
спецификации и исходного кода не противоречит действующему
законодательству и нормативным документам Российской Федерации. Напротив,
её применение соответствует современным требованиям к надёжности,
корректности и безопасности программного обеспечения, а также может
рассматриваться как элемент более общего процесса безопасной разработки и
контроля качества программных систем
\cite{constitution_rf,fz149_info,doctrine_ib_rf,gost_56939_2024,fstek_240_2023}.
